\section{最终投资结论}
本报告的最终建议不是“一篮子买光通信”，而是\textbf{只买兑现、限制期权、承认情绪、等待验证}。截至 2026-06-26 收盘价，覆盖组合的内在价值锚加权空间为 -47.8\%，纳入市场隐含预期与券商锚后的综合空间为 -23.1\%。这意味着板块仍有强产业逻辑，但从组合角度看已经不是低估扩散阶段，而是业绩兑现和市场情绪共同定价阶段。投资动作应从主题追涨切换为三件事：第一，核心模块只在利润、毛利率和现金流继续兑现时加仓；第二，光纤光缆、网络设备和通信线缆只按慢周期现金流和项目节奏持有观察；第三，上游芯片/器件和设备材料只给期权仓位，不能用长期空间掩盖短期分母不足。

\section{建议的研究组合结构}
以下权重是研究组合权重，不是账户交易指令。它用于约束不同产业链层级的相对重要性，防止因为光通信主题强就无差别追涨。当前基准情景下，组合应维持\textbf{中性偏进攻观察}：保留核心模块敞口，但新增资金必须等待 Q2/Q3 交付、毛利率和现金流验证。

\begin{exhibitbox}[表：最终组合建议矩阵]
\centering
\small
\begin{tabularx}{\exhibitboxwidth}{>{\bfseries\raggedright\arraybackslash}p{2.1cm} >{\centering\arraybackslash}p{1.4cm} >{\raggedright\arraybackslash}p{2.8cm} >{\raggedright\arraybackslash\sloppy\hspace{0pt}}X >{\raggedright\arraybackslash\sloppy\hspace{0pt}}p{3.0cm}}
\toprule
\textbf{篮子} & \textbf{研究权重} & \textbf{代表标的} & \textbf{当前建议} & \textbf{升/降级触发} \\
\midrule
核心兑现池 & 30--35\% & 中际旭创、新易盛 & 作为组合主线，但不无条件追高；利润、毛利率、现金流三项确认后才提高权重。 & Q2/Q3 净利环比继续改善且现金流不背离，上调；若毛利率或订单节奏破坏，下调。 \\
市场支持观察池 & 25--30\% & 中天科技、长飞光纤、亨通光电、通鼎互联、永鼎股份 & 承认市场情绪锚，但只按慢周期现金流和项目节奏管理，不能套用纯 AI 模块倍数。 & 光纤光缆价格、海缆/运营商项目和经营现金流改善，上调；成交额和现金流同时回落，下调。 \\
上游战略期权池 & 10--15\% & 源杰科技、长光华芯、仕佳光子、光库科技、腾景科技 & 只给期权仓位，重点看认证、良率、客户数和收入占比。 & 客户认证转量产且毛利率稳定，上调；只有题材无收入，维持低配。 \\
网络设备与应用层 & 10--15\% & 中兴通讯、烽火通信 & 稳健观察，按运营商 capex、云网络订单和利润率定价。 & 运营商/云网络项目和利润率改善，上调；非光通信业务稀释或项目延后，下调。 \\
低配/等待修复池 & 0--10\% & 光迅、华工、剑桥、太辰光、德科立、联特、兆龙、意华、神宇及设备材料观察池 & 只有在目标价空间、利润分母或现金流修复后才提高权重；亏损或近零 EPS 标的不强行买入。 & 综合空间转正且证据质量提升，上调；预期继续前置而交付不足，维持低配或减持。 \\
\bottomrule
\end{tabularx}
\end{exhibitbox}

\section{重点标的执行清单}
最终建议必须落到单票执行层。下表把当前价、综合目标、空间、动作和下一步触发绑定在一起。中天科技、长飞光纤、通鼎互联这类标的的重点不是机械减持，而是承认市场已经给出情绪溢价，同时要求后续现金流、订单和项目节奏确认；若确认失败，综合目标会向内在价值锚收敛。

\begin{exhibitbox}[表：重点标的最终执行清单]
\centering
\tiny
\begin{tabularx}{\exhibitboxwidth}{>{\bfseries\raggedright\arraybackslash}p{1.7cm} >{\raggedleft\arraybackslash}p{0.85cm} >{\raggedleft\arraybackslash}p{0.95cm} >{\raggedleft\arraybackslash}p{0.85cm} >{\centering\arraybackslash}p{1.25cm} >{\raggedright\arraybackslash\sloppy\hspace{0pt}}X >{\raggedright\arraybackslash\sloppy\hspace{0pt}}X}
\toprule
\textbf{标的} & \textbf{现价} & \textbf{目标} & \textbf{空间} & \textbf{动作} & \textbf{加仓/上修条件} & \textbf{降级/退出条件} \\
\midrule
中际旭创 300308 & 1253.89 & 1229.56 & -1.9\% & \stance{riskamber}{中性} & Q2/Q3 净利、毛利率、经营现金流同步验证，1.6T 客户交付没有延迟；回撤到目标区间下沿附近但基本面未破时优先补。 & 净利环比低于 Q1、毛利率下滑或大客户订单节奏放缓，降为只观察不加仓。 \\
新易盛 300502 & 566.00 & 576.55 & +1.9\% & \stance{riskamber}{中性} & 高增长订单继续兑现，单一客户风险没有扩大，现金流跟上利润；若综合空间重新扩大到 8\% 以上，进入增持候选。 & 高增长只体现在收入不体现在现金流，或 2027E 增速预期下修，降低权重。 \\
中天科技 600522 & 61.29 & 47.81 & -22.0\% & \stance{riskamber}{中性观察} & 海缆、电网、通信线缆订单和现金流确认；市场情绪锚维持但不能脱离经营现金流。 & 光纤价格下行、海缆/电网项目延后，或成交额回落导致情绪锚失效。 \\
长飞光纤 601869 & 543.00 & 434.40 & -20.0\% & \stance{riskamber}{中性观察} & 预制棒/光纤/光缆利用率回升，数据中心需求真正进入收入和现金流；只做强共识观察，不追高。 & 利用率、毛利率或经营现金流未改善，强共识溢价回落时下调市场锚。 \\
通鼎互联 002491 & 36.15 & 26.03 & -28.0\% & \stance{riskamber}{中性观察} & 通信线缆和项目交付带动利润分母修复，Q2/Q3 经营现金流改善；按慢周期底盘处理。 & 利润改善不能延续，或 AI 相关纯度被证伪，退出核心观察。 \\
源杰科技 688498 & 1776.91 & 1421.53 & -20.0\% & \stance{riskamber}{中性观察} & 激光芯片客户认证、收入规模和毛利率同时提升；只能作为上游战略期权，不能按模块龙头仓位处理。 & 认证不及预期、收入放量慢于股价隐含预期，维持中性观察或降级。 \\
光迅科技 002281 & 240.34 & 187.47 & -22.0\% & \stance{riskamber}{中性观察} & 光芯片/器件和模块业务利润率修复，证明平台价值能穿越板块情绪。 & 估值继续高于利润兑现速度，或现金流弱于利润，维持中性观察。 \\
中兴通讯 000063 & 35.67 & 31.06 & -12.9\% & \stance{riskamber}{中性} & 运营商、云网络和服务器相关业务利润率改善，现金流维持稳健；作为应用层稳健观察。 & 非光通信业务稀释加重，运营商 capex 延后，维持中性或降低权重。 \\
光库科技 300620 & 362.43 & 253.70 & -30.0\% & \stance{riskamber}{中性观察} & 薄膜铌酸锂/调制器期权进入订单和收入，且综合空间不再深负。 & 股价继续把长期期权前置而收入和利润不跟，维持减持。 \\
联特科技 301205 & 315.86 & 116.17 & -63.2\% & \stance{riskred}{减持} & 高 beta 模块订单、客户认证和利润分母修复；未修复前不因题材追涨。 & Q2/Q3 利润仍薄或现金流弱，维持减持。 \\
\bottomrule
\end{tabularx}
\end{exhibitbox}

\section{三种情景下的操作剧本}
\begin{exhibitbox}[表：情景、动作与复盘纪律]
\centering
\small
\begin{tabularx}{\exhibitboxwidth}{>{\bfseries\raggedright\arraybackslash}p{1.9cm} >{\raggedright\arraybackslash\sloppy\hspace{0pt}}X >{\raggedright\arraybackslash\sloppy\hspace{0pt}}X >{\raggedright\arraybackslash\sloppy\hspace{0pt}}X}
\toprule
\textbf{情景} & \textbf{触发条件} & \textbf{组合动作} & \textbf{重点复盘} \\
\midrule
牛市兑现 & Q2/Q3 模块龙头净利环比强增长，1.6T 订单批量化，上游认证转收入，光纤光缆现金流改善。 & 提高核心模块权重，选择性增加上游战略期权和慢周期底盘；仍不扩大亏损/近零 EPS 设备材料敞口。 & 毛利率是否稳定、经营现金流是否跟上、客户集中度是否恶化。 \\
基准验证 & 模块利润继续增长但估值已经充分，慢周期链条只有部分项目兑现，上游认证仍在推进。 & 维持中性偏进攻观察；核心池逢基本面确认加，市场支持观察池不追涨，低配分母不足标的。 & 现价是否重新低于综合目标区间下沿，现金流和订单能否支持情绪锚。 \\
预期破裂 & Q2/Q3 利润低于 Q1，毛利率或 ASP 下行，经营现金流弱于利润，光纤光缆/网络项目延后。 & 降低主题总敞口，优先减掉期权和低证据质量标的；中性观察标的若情绪锚失效，降为减持。 & 内在价值锚下修幅度、成交额分位回落、券商目标和一致预期是否同步下修。 \\
\bottomrule
\end{tabularx}
\end{exhibitbox}

\section{复盘纪律与结论}
未来一个季度的复盘顺序应固定为：先看利润分母，再看现金流，再看订单/认证，最后才看市场情绪。若某只股票只有成交额和题材热度，没有利润、毛利率或现金流确认，不能因为产业链位置好就上调动作；若某只股票内在价值偏低但市场锚、成交和项目兑现仍强，也不能机械写成减持，而应维持中性观察并列出情绪溢价失效条件。最终结论是：光通信仍是 AI 硬件链最重要方向之一，但当前阶段的正确做法不是扩大无差别敞口，而是在中际旭创、新易盛等利润已兑现标的上等待确认，在中天科技、长飞光纤、通鼎互联等市场支持观察标的上用现金流验证情绪锚，在上游芯片/器件和设备材料上保持期权仓位和严格止错。

\begin{exhibitbox}[表：下一次更新检查清单]
\centering
\small
\begin{tabularx}{\exhibitboxwidth}{>{\bfseries\raggedright\arraybackslash}p{2.4cm} >{\raggedright\arraybackslash\sloppy\hspace{0pt}}X >{\raggedright\arraybackslash\sloppy\hspace{0pt}}X}
\toprule
\textbf{检查项} & \textbf{上修条件} & \textbf{下修条件} \\
\midrule
利润分母 & Q2/Q3 净利润继续环比改善，且增长不只来自一次性收益。 & 净利润低于 Q1 或低于市场隐含增速，先下调 EPS 和目标价。 \\
毛利率与 ASP & 模块毛利率稳定，上游器件和芯片放量不牺牲毛利率。 & ASP 下行、毛利率下降或客户议价增强，降低估值倍数。 \\
经营现金流 & 经营现金流跟上净利润，应收和存货没有快于收入恶化。 & 现金流持续弱于利润，说明订单质量低于股价隐含预期。 \\
订单与认证 & 1.6T、CPO、硅光、激光芯片、光纤光缆项目进入收入确认。 & 只停留在送样、认证或主题宣传，不能上调动作标签。 \\
市场情绪锚 & 成交额分位、券商目标和项目兑现同时支撑市场锚。 & 成交回落、券商目标下修或项目延后，情绪锚向内在价值锚收敛。 \\
\bottomrule
\end{tabularx}
\end{exhibitbox}
